\begin{solapartat}
En el balanç energètic de l'efecte fotoelèctric tenim:
\[ h\frac{c}{\lambda} - W = E_C \ \text{\punts{0,4}} \Rightarrow \]
\begin{align*}
W = h\frac{c}{\lambda} - E_c &= 6,63\cdot 10^{-34}\,\frac{3\cdot 10^{8}}{23,7\cdot 10^{-9}} - 47,7\,\frac{1,6\cdot 10^{-19}}{1\,\mathrm{eV}}\\
&= 7,60\cdot 10^{-19}\ \mathrm{J} \ \text{\punts{0,3}} = 4,75\ \mathrm{eV} \ \text{\punts{0,3}}
\end{align*}
\end{solapartat}

\begin{solapartat}
La longitud d'ona llindar la obtindrem fent que l'energia cinètica dels electrons emesos sigui zero. \punts{0,2}
\[ \lambda_L = h\frac{c}{W} = 2,62\cdot 10^{-7}\ \mathrm{m} \ \text{\punts{0,5}} \]
La longitud d'ona llindar no depèn de la potencia de la radiació incident, per tant si dupliquem aquesta potència la longitud d'ona llindar no variarà \punts{0,3}
\end{solapartat}
